\documentclass[10pt,a4paper]{amsart}
\usepackage[margin=19mm]{geometry}
\usepackage{amsmath,amssymb,mathtools}
\usepackage[scr=rsfs,cal=euler]{mathalfa}
\usepackage{xfrac}
\usepackage[unicode,hidelinks,bookmarks=false]{hyperref}
\newcommand{\ie}{i.e.\ }
\newcommand{\cf}{cf.\ }
\newcommand{\resp}{respectively\ }
\newcommand{\Subsection}{\subsection}
\providecommand{\nobreakdash}{\nobreak\hbox{-}}
\setlength{\emergencystretch}{2em}
\multlinegap0pt
\usepackage{bm}
\usepackage{bbm}
\usepackage{dsfont}
\usepackage{xspace}
\usepackage{cancel}
\usepackage{mathtools}
\usepackage{amsthm}
\makeatletter
\@ifundefined{theorem}{\newtheorem{theorem}{Theorem}}{}
\@ifundefined{lemma}{\newtheorem{lemma}[theorem]{Lemma}}{}
\@ifundefined{proposition}{\newtheorem{proposition}[theorem]{Proposition}}{}
\makeatother
\title[Thermodynamic Formalism for a Class of Hyperbolic Transcende]{Thermodynamic Formalism for a Class of Hyperbolic Transcendental Meromorphic Functions}
\author{Hamid Naderiyan}
\date{Source: arXiv:2506.06760v1 (CC BY 4.0)}
\begin{document}
\maketitle

\noindent\emph{Source and licence.} Thermodynamic Formalism for a Class of Hyperbolic Transcendental Meromorphic Functions. Version arXiv:2506.06760v1 (CC BY 4.0), distributed under the Creative Commons Attribution 4.0 International licence, as recorded in the arXiv licence metadata for this exact version and shown on the abstract page https://arxiv.org/abs/2506.06760v1. Full rights notice: licenses/arxiv-2506.06760-CC-BY-4.0.txt.\par\smallskip

\noindent\textit{Editorial scope.} Counted unit: the Lemma whose source label is \texttt{quasi-invariant} in the article (source0.tex lines 605--617, source label \texttt{quasi-invariant}), delivered together with the article's own complete original proof of it (source0.tex lines 618--643, one \texttt{proof} environment of 26 lines). No mathematics is written, repaired or re-derived: statement and proof are the article's own text line for line. Required context retained uncounted: delta (lines 182--184); T (lines 208--210); estimate for tau derivative stronger (lines 250--253); Borel theorem 2 (lines 301--303); exp Ionescu S\_n (lines 369--376); equivalence of eigenmeasure and conformal measure (lines 426--436). Every definition, display and label of those lines is retained unchanged. Omitted from this package and not redistributed: 1--181, 185--207, 211--249, 254--300, 304--368, 377--425, 437--604, 644--778. Nothing inside a retained excerpt is omitted or altered: every retained body is delivered exactly as the article's lines are written, so no omission receipt of any kind accompanies this package. Citations: the retained excerpts carry no citation command at all, so no bibliography entry of the article is transcribed and no reference is redistributed. Reference adaptation: none was needed and none was made; every reference inside the delivered unit resolves inside it, so no unresolved reference or citation is produced. Printed slips kept literal: no printed word, symbol, hypothesis or number of the counted statement or of its proof was altered. Layout and class: the article's own class is \texttt{amsart} with packages including inputenc, babel, amsmath, amsfonts, amssymb, amsthm, siunitx, bbm, hyperref, natbib, comment, tasks, enumitem, stix; the wrapper is the lane's pinned \texttt{amsart} editorial wrapper at 10pt on A4, which restates the theorem-like environments the retained text needs and adds the article's own macro definitions that the retained text needs (none), with extra wrapper packages (bm, bbm, dsfont, xspace, cancel, mathtools). The delivered numbering is the unit's own. Assets: no retained span contains a figure environment, a graphics-inclusion command, a PDF-page inclusion, a table or a TikZ picture; no image file is copied into this package, the package declares no asset dependency and no image file is redistributed with it. Licence: version arXiv:2506.06760v1 of this article is distributed under the Creative Commons Attribution 4.0 International licence, as recorded in the arXiv licence metadata for this exact version and shown on the abstract page https://arxiv.org/abs/2506.06760v1. A full rights notice is delivered at licenses/arxiv-2506.06760-CC-BY-4.0.txt. Characters: every retained excerpt is pure ASCII, so the delivered engine, XeLaTeX with its default Latin Modern text font, reports no missing character.
\begin{equation}\label{delta}
    \delta:=\frac{1}{4}\textup{dist}(\mathcal{J}(f),\overline{\mathcal{P}}_f)
\end{equation}

\begin{equation}\label{T}
    |z|, |f(z)| \geq T>0, \hspace{.5cm} z \in\mathcal{J}(f)
\end{equation}

\begin{lemma}\label{estimate for tau derivative stronger}
    There exists $c>0$ such that for every $w \in \mathcal{J}(f)$ and every $z\in f^{-1}(\{w\})$ we have,
    $$\left|(f_z^{-1})'(w)\right|_{\tau}^{t}=|f'(z)|_{\tau}^{-t} \leq c^{t}\dfrac{|w|^{-(1+\frac{1}{M}-\tau)t}}{|z|^{(\tau-1)t}}.$$
\end{lemma}

\begin{equation}\label{Borel theorem 2}
    \sum_{f(z)=w}\dfrac{1}{|z|^{(\tau-1)t}} \leq M_{(\tau-1)t}.
\end{equation}

    \begin{lemma}\label{exp Ionescu S_n}
        There exists $K>1$ such that for every $n \in \mathbb{N}$, every $w_1,w_2 \in \mathcal{J}(f)$ with $|w_1-w_2|<\delta$ and  
        every $z \in f^{-n}(\{w_1 \})$, we have
        \begin{align*}
            \left| \exp\left(S_n\Phi_t\left(f^{-n}_z(w_1)\right)\right) -\exp\left(S_n\Phi_t\left(f^{-n}_z(w_2)\right)\right) \right| &= \left|\left|(f_z^{-n})'(w_1)\right|_{\tau}^t - \left|(f_z^{-n})'(w_2)\right|_{\tau}^t\right|\\
            & \leq te^{tK}K\left|(f_z^{-n})'(w_1)\right|_{\tau}^t|w_1-w_2|.
        \end{align*}
    \end{lemma}

\begin{proposition}\label{equivalence of eigenmeasure and conformal measure}
     Given a constant $c$ and a $ce^{-\Phi_t}-$conformal measure $m$ on $\mathcal{J}(f)$. The following conditions are equivalent,
     
    (a) $m(f^n(A))=c^n\int_A\exp(-S_n\Phi_t)\text{d}m$, for every $n\in \mathbb{N}$ and every Borel $A\subseteq \mathcal{J}(f)$ where $f^n$ is injective on $A$.

    (b) $m(f(A))=c\int_A\exp(-\Phi_t)\text{d}m$, for every Borel $A\subseteq \mathcal{J}(f)$ where $f$ is injective on $A$.

    (c) $\left(\mathcal{L}_t \right)^*m=cm$

    (d) $\left(\mathcal{L}_t^n \right)^*m=c^nm$ for every $n\in \mathbb{N}$.
\end{proposition}

\begin{lemma}\label{quasi-invariant}
    Given a constant $c$ and a $ce^{-\Phi_t}-$conformal measure $m$ for $f$. Then the following items hold.
    \begin{itemize}
        \item[(a)] For every $R>0$, there exists $c_R>0$ such that for every Borel set  $B \subseteq \mathcal{J}(f) \cap B(R)$,
    $$m(f^{-1}(B))\leq c_R m(B),$$
    where $c_R \to 0$, as $R\to \infty$.
        \item[(b)]  $m$ is quasi-invariant, i.e. $m(f^{-1}(B))=0$ when $m(B)=0$.
        \item[(c)] There exist $R>0$ and $0<c_R<1$ such that for every $n \in \mathbb{N}$ and every Borel set  $B \subseteq \mathcal{J}(f) \cap B(R)$,
        $$m \left( \cap_{i=0}^{n}f^{-i}(B) \right) \leq c_R^nm(B).$$
        \item[(d)]  For every $w_1\in \mathcal{J}(f)$, every $z\in f^{-n}(\{w_1\})$ and every Borel set  $B \subseteq \mathcal{J}(f) \cap D(w_1,\delta)$,
    $$m(f_z^{-n}(B))= c^{-n}\int_B \exp\left(S_n\Phi_t(f_z^{-n}(w))\right)dm(w).$$
    \end{itemize}
\end{lemma}

\begin{proof}
     Given $f(z)=w\in \mathcal{J}(f)$ and a Borel set $B\subseteq D(w,\delta) \cap \mathcal{J}(f)$, where $\delta$ was given in the equation (\ref{delta}). By lemma \ref{exp Ionescu S_n} one can find $K_t>1$ such that for all $w_1\in B,$
     $$\exp\left(\Phi_t\left(f^{-1}(w_1)\right)\right)=\left|(f_z^{-1})'(w_1)\right|_{\tau}^{t}\leq K_t \left|(f_z^{-1})'(w)\right|_{\tau}^{t}.$$
     Therefore, by lemma \ref{estimate for tau derivative stronger} and the fact that $f: f^{-1}_z(B) \to B$ is bijective we obtain,
     $$m(B)=m(f(f^{-1}_z(B)))=\int_{f_z^{-1}(B)}ce^{-\Phi_t}\text{d}m \geq cK_t^{-1}(d^{t})^{-1}\dfrac{|w|^{(1+\frac{1}{M}-\tau)t}}{|z|^{-(\tau-1)t}} m(f_z^{-1}(B)).$$
     This means there exists $b_t>0$ such that 
     $$m(f_z^{-1}(B)) \leq b_t \dfrac{|w|^{-(1+\frac{1}{M}-\tau)t}}{|z|^{(\tau-1)t}} m(B).$$
     and therefore by the estimate (\ref{Borel theorem 2}) we have,
     \begin{align*}
         m(f^{-1}(B))=\sum_{f(z)=w}m(f_z^{-1}(B)) &\leq b_tm(B)|w|^{-(1+\frac{1}{M}-\tau)t}\sum_{f(z)=w}\dfrac{1}{|z|^{(\tau-1)t}}\\
         &\leq  b_tm(B)|w|^{-(1+\frac{1}{M}-\tau)t} M_{(\tau-1)t}
     \end{align*}
     Given $R>0$ and a Borel set $B\subseteq \mathcal{J}(f) \cap B(R)$, one can find a countable set $\{w_i\}_i \subseteq B$ such that 
     $B \subseteq \cup_i D(w_i,\delta)$. Inductively, we can obtain $B_i\subseteq D(w_i,\delta)$ such that $B=\cup_i B_i$ where $B_{i}$'s are disjoint. Then the above estimate gives,
     \begin{align*}
    m\left(f^{-1}(B)\right) = \sum_i m\left(f^{-1}(B_{i})\right) &\leq \sum_i b_t m(B_i)|w_{i}|^{-(1+\frac{1}{M}-\tau)t}M_{(\tau-1)t}\\
    & \leq \sum_i b_tm(B_i)R^{-(1+\frac{1}{M}-\tau)t}M_{(\tau-1)t}\\
    & \leq b_tR^{-(1+\frac{1}{M}-\tau)t}M_{(\tau-1)t}m(B).
     \end{align*}
     This proves item (a).\\
     Item (b) follows from item (a) with regards to (\ref{T}).\\
     For item (c), first one can choose $R>0$ large enough such that $c_R<1$ by item (a). Then $m(B\cap f^{-1}(B))\leq m(f^{-1}(B))\leq c_Rm(B)$. Therefore, by the following estimate
     $$m\left( \cap_{i=0}^nf^{-i}(B)\right) \leq m\left( \cap_{i=1}^nf^{-i}(B)\right)=m\left( f^{-1}\left(\cap_{i=0}^{n-1}f^{-i}(B)\right)\right),$$
     one can use induction to establish item (c).\\
     Finally, item (d) follows from proposition \ref{equivalence of eigenmeasure and conformal measure} item (d).
\end{proof}

\end{document}
